% Target: rmp-iii-a-torus-two
% Formula: A five-tensor enlarged MPO word, with its fourth tensor marked.
% Ink: Kernel flat frame; the periodic word is a closed operator string behind
%   the boxes and the marked tensor uses the marked semantic role.
% Boundary: The virtual MPO bond is periodic; all ten physical legs are open.
% Source: paper TN-Review-main.tex line 1594; author source ImagesReview Section
%   III/ImagesReview Section III.tex lines 553-562; archival output
%   fig3_fig3-pepe_Eq62.pdf
% Capabilities: kernel, mpo-word, periodic-closure, marked-insertion
\[
  \tenkzkernel
  \begin{tenkz}[rows={wire}, cols=5, bonds=none, physical=updown]
    \tn[at=(1,1), name=b1, skin=box]{}
    \tn[at=(1,2), name=b2, skin=box]{}
    \tn[at=(1,3), name=b3, skin=box]{}
    \tn[at=(1,4), name=b4, skin=box, role=marked]{}
    \tn[at=(1,5), name=b5, skin=box]{}
    \tnwire[kind=string, species=operator]{b1.e}{b2.w}
    \tnwire[kind=string, species=operator]{b2.e}{b3.w}
    \tnwire[kind=string, species=operator]{b3.e}{b4.w}
    \tnwire[kind=string, species=operator]{b4.e}{b5.w}
    % the periodic return of the word, drawn as the source's long loop below,
    % as a junction chain of invisible corner atoms (via= retired, #6897)
    \tn[at=0.8 e of (1,5), skin=none, name=jwrap1]{}
    \tn[at=1.1 s of 0.8 e of (1,5), skin=none, name=jwrap2]{}
    \tn[at=1.1 s of 0.8 w of (1,1), skin=none, name=jwrap3]{}
    \tn[at=0.8 w of (1,1), skin=none, name=jwrap4]{}
    \tnwire[kind=string, species=operator, name=wrap1]{b5.e}{jwrap1}
    \tnwire[kind=string, species=operator, name=wrap2,
            cross={over at crossing of wrap1 and wrap2}]{jwrap1}{jwrap2}
    \tnwire[kind=string, species=operator, name=wrap3,
            cross={over at crossing of wrap2 and wrap3}]{jwrap2}{jwrap3}
    \tnwire[kind=string, species=operator, name=wrap4,
            cross={over at crossing of wrap3 and wrap4}]{jwrap3}{jwrap4}
    \tnwire[kind=string, species=operator, name=wrap5,
            cross={over at crossing of wrap4 and wrap5}]{jwrap4}{b1.w}
  \end{tenkz}
\]
